\documentclass[11pt,a4paper]{article}
\usepackage[T2A]{fontenc}
\usepackage[utf8]{inputenc}
\usepackage[russian]{babel}
\usepackage{geometry}
\usepackage{amsmath,amssymb,mathtools}
\usepackage{array,booktabs}
\usepackage{xcolor}
\usepackage{hyperref}
\usepackage{tikz}
\geometry{left=22mm,right=22mm,top=22mm,bottom=23mm}
\hypersetup{unicode=true,colorlinks=true,linkcolor=blue!55!black,
pdftitle={Ревью домашней работы},pdfauthor={Лёвин Артём Александрович}}
\setlength{\parindent}{0pt}
\setlength{\parskip}{0.40em}
\setlength{\emergencystretch}{2em}
\renewcommand{\arraystretch}{1.4}
\begin{document}
\begin{titlepage}
\thispagestyle{empty}
\vspace*{48mm}
\begin{center}
{\Huge\bfseries Ревью домашней работы\par}
\vspace{23mm}
{\Large Автор: Лёвин Артём Александрович\par}
\vspace{10mm}
{\large 9 октября 2026 года\par}
\end{center}
\vfill
\begin{center}
\begin{tikzpicture}
\draw[blue!30!black,line width=0.65pt]
(0,0) .. controls (1.7,0.58) and (3.3,-0.58) .. (5,0)
.. controls (6.6,0.58) and (8.2,-0.58) .. (9.9,0);
\end{tikzpicture}
\end{center}
\vspace{8mm}
\end{titlepage}

\clearpage
\section*{Сводная таблица}
\small
\begin{center}
\begin{tabular}{@{}>{\raggedright\arraybackslash}p{0.9cm}>{\raggedright\arraybackslash}p{2.9cm}>{\raggedright\arraybackslash}p{4.5cm}>{\raggedright\arraybackslash}p{2.1cm}>{\raggedright\arraybackslash}p{3.2cm}@{}}
\toprule
\textbf{№} & \textbf{Тема} & \textbf{Суть ошибки} & \textbf{Ваш ответ} & \textbf{Правильный ответ}\\
\midrule
\hyperref[task:3]{3} & Степени с иррациональными показателями & При переходе к основанию $\frac12$ потеряно слагаемое $2$ в показателе степени знаменателя. & $2$ & $\frac12$\\
\addlinespace[6pt]
\hyperref[task:10]{10} & Кубические корни и степени & Решение оборвано на записи дроби; показано выражение для числителя, но итоговая дробь не получена. & Не указан & $x^2\sqrt[3]{x^2}$\\
\bottomrule
\end{tabular}
\end{center}
\normalsize

\clearpage
\vspace*{\fill}
\begin{center}
\section*{Разбор ошибок}
\end{center}
\vspace*{\fill}
\clearpage

\phantomsection\label{task:3}
\section*{Задание № 3}
\subsection*{Текст задания}
Упростите выражение:
\[
\frac{(0{,}5)^{\sqrt{10}-1}}{2^{\,2-\sqrt{10}}}.
\]
\subsection*{Ваше решение}
В записи после перехода к основанию $\frac12$ знаменатель представлен как
$\bigl(\frac12\bigr)^{\sqrt{10}}$. После этого Вы вычитаете показатели
степеней и получаете:
\[
\frac{\left(\frac12\right)^{\sqrt{10}-1}}
{\left(\frac12\right)^{\sqrt{10}}}
=\left(\frac12\right)^{\sqrt{10}-1-\sqrt{10}}
=\left(\frac12\right)^{-1}=2.
\]
Это верное вычисление для \emph{записанной после преобразования дроби},
но сама эта дробь не равна исходному выражению.
\subsection*{Ваш ответ}
\[\boxed{2}.\]
\subsection*{Ошибка}
\textcolor{red}{Ошибка:} при переписывании знаменателя потеряно слагаемое
$2$ в показателе степени. Поэтому основание было изменено с $2$ на
$\frac12$ неверным образом.
\subsection*{Где именно возникает ошибка}
\textbf{Последняя правильная запись:}
\[
\frac{(0{,}5)^{\sqrt{10}-1}}{2^{2-\sqrt{10}}}.
\]
\textbf{Первый неверный переход:}
\[
\underbrace{2^{2-\sqrt{10}}}_{\text{исходный знаменатель}}
\ \color{red}{\ne}\
\underbrace{\left(\frac12\right)^{\sqrt{10}}}_{\text{знаменатель в Вашем преобразовании}}.
\]
Действительно,
\[
2^{2-\sqrt{10}}=2^2\cdot 2^{-\sqrt{10}}
=4\cdot\left(\frac12\right)^{\sqrt{10}},
\]
а не просто $\left(\frac12\right)^{\sqrt{10}}$.
\subsection*{Почему этот переход неверен}
При переходе от основания $2$ к обратному основанию $\frac12$
\emph{весь показатель степени} меняет знак:
\[
2^t=\left(\frac12\right)^{-t}.
\]
Если $t=2-\sqrt{10}$, то
\[
-t=-(2-\sqrt{10})=\sqrt{10}-2.
\]
Следовательно,
\[
\boxed{2^{2-\sqrt{10}}
=\left(\frac12\right)^{\sqrt{10}-2}}.
\]
Число $2$, стоявшее в исходном показателе, нельзя опускать:
от него зависит множитель $2^2=4$ в знаменателе.

Именно поэтому полученный Вами ответ оказался в четыре раза больше
правильного. При исключении множителя $4$ из знаменателя значение всей
дроби увеличивается в четыре раза: вместо $\frac12$ получается $2$.
\subsection*{Ключевая идея}
\textcolor{blue!60!black}{Ключевая идея:} при действиях со степенями
удобно сначала записать числитель и знаменатель как степени одного
основания, \emph{сохранив каждый член показателя}, а затем воспользоваться
правилом
\[
\frac{a^p}{a^q}=a^{p-q},\qquad a>0.
\]
\subsection*{Исправление от последнего правильного шага}
Сохраним первоначальный знаменатель и заменим только десятичную дробь
$0{,}5$ на $2^{-1}$:
\[
\frac{(0{,}5)^{\sqrt{10}-1}}{2^{2-\sqrt{10}}}
=\frac{(2^{-1})^{\sqrt{10}-1}}{2^{2-\sqrt{10}}}
=\frac{2^{1-\sqrt{10}}}{2^{2-\sqrt{10}}}.
\]
Теперь основания одинаковы, и при делении показатели вычитаются:
\[
\frac{2^{1-\sqrt{10}}}{2^{2-\sqrt{10}}}
=2^{(1-\sqrt{10})-(2-\sqrt{10})}
=2^{-1}=\frac12.
\]
\subsection*{Полное правильное решение}
\textbf{Шаг 1.} Перепишем десятичную дробь:
\[
0{,}5=\frac12=2^{-1}.
\]
\textbf{Шаг 2.} Применим правило степени степени
$(a^u)^v=a^{uv}$:
\[
(0{,}5)^{\sqrt{10}-1}
=(2^{-1})^{\sqrt{10}-1}
=2^{-(\sqrt{10}-1)}
=2^{1-\sqrt{10}}.
\]
\textbf{Шаг 3.} Подставим полученное выражение в числитель:
\[
\frac{(0{,}5)^{\sqrt{10}-1}}{2^{2-\sqrt{10}}}
=\frac{2^{1-\sqrt{10}}}{2^{2-\sqrt{10}}}.
\]
\textbf{Шаг 4.} Вычтем показатель знаменателя из показателя
числителя, не забывая раскрыть скобки:
\begin{align*}
(1-\sqrt{10})-(2-\sqrt{10})
 &=1-\sqrt{10}-2+\sqrt{10}\\
 &=-1.
\end{align*}
\textbf{Шаг 5.} Используем определение отрицательной степени:
\[
2^{-1}=\frac1{2^1}=\frac12.
\]
Значит, всё исходное выражение равно $\frac12$.
\subsection*{Проверка}
Умножим найденное значение на исходный знаменатель.
Должен получиться числитель:
\[
\frac12\cdot 2^{2-\sqrt{10}}
=2^{-1+2-\sqrt{10}}
=2^{1-\sqrt{10}}
=(0{,}5)^{\sqrt{10}-1}.
\]
Числитель восстановлен точно; проверка выполнена.
\subsection*{Итог}
\textcolor{teal!60!black}{Итог:} ошибка связана не с вычитанием
показателей, а с потерей слагаемого при смене основания степени.
Правильный ответ:
\[
\boxed{\frac12}.
\]
\subsection*{Задача на закрепление}
Упростите выражение:
\[
\frac{(0{,}5)^{\sqrt7-2}}{2^{\,4-\sqrt7}}.
\]

\clearpage
\phantomsection\label{task:10}
\section*{Задание № 10}
\subsection*{Текст задания}
При $x\ne0$ упростите выражение:
\[
\frac{\bigl(\sqrt[3]{7x^6}\bigr)^2}{\sqrt[3]{49x^4}}.
\]
\subsection*{Ваше решение}
После знака равенства в работе начата новая дробь. Под дробной
чертой видно $7^{2/3}x^4$, но запись над чертой и дальнейшие
преобразования отсутствуют.

Само выражение $7^{2/3}x^4$ может быть получено при правильном
упрощении \emph{числителя} исходной дроби:
\[
\bigl(\sqrt[3]{7x^6}\bigr)^2
=\bigl(7^{1/3}x^2\bigr)^2
=7^{2/3}x^4.
\]
Однако в Вашей записи эта часть находится под дробной чертой.
Поэтому нельзя считать, что итоговая дробь записана корректно,
или додумывать отсутствующие шаги.
\subsection*{Ваш ответ}
Окончательный ответ не указан.
\subsection*{Ошибка}
\textcolor{red}{Ошибка:} решение не завершено. После исходной дроби
появилась неполная запись; выражение, соответствующее преобразованию
числителя, помещено под дробную черту. Знаменатель и результат
деления не вычислены.
\subsection*{Где именно возникает ошибка}
\textbf{Последний полностью правильный шаг:} исходная запись задания.
Отдельно найденный множитель $7^{2/3}x^4$ действительно равен
\emph{числителю}, но не знаменателю.

\textbf{Первый неполный шаг:} после знака равенства не записано,
что стоит в числителе новой дроби, а выражение под чертой
не является упрощённым знаменателем исходной дроби.
Чтобы проверить это, отдельно рассмотрим знаменатель:
\[
\sqrt[3]{49x^4}
=\sqrt[3]{7^2}\,\sqrt[3]{x^4}
=7^{2/3}\sqrt[3]{x^4}.
\]
Здесь \emph{кубический корень из $x^4$ сохраняется}.
В общем случае
\[
\sqrt[3]{x^4}\ne x^4.
\]
Например, при $x=2$ получаются $\sqrt[3]{16}$ и $16$ ---
это разные числа.
\subsection*{Почему это неверно}
Дробная черта означает деление числителя на знаменатель.
Изменение положения выражения относительно дробной черты меняет
его роль в вычислении: то, что находится сверху, делится на то,
что находится снизу.

Кроме того, у корней и степеней есть два различных правила:
\[
\bigl(\sqrt[3]{a}\bigr)^2=\sqrt[3]{a^2},
\qquad
\sqrt[3]{a^4}=\sqrt[3]{a^3\cdot a}
=a\sqrt[3]{a}.
\]
Второе равенство верно и для отрицательных действительных $a$,
поскольку речь идёт о кубическом корне.
Поэтому нельзя просто убрать знак кубического корня у $x^4$.

Также необходимо учитывать условие $x\ne0$: оно гарантирует,
что исходный знаменатель $\sqrt[3]{49x^4}$ не равен нулю
и деление допустимо.
\subsection*{Ключевая идея}
\textcolor{blue!60!black}{Ключевая идея:} здесь удобнее всего
объединить кубические корни. Возведение кубического корня в квадрат
можно перенести под знак корня:
\[
\bigl(\sqrt[3]{A}\bigr)^2=\sqrt[3]{A^2}.
\]
Затем кубические корни числителя и знаменателя объединяются:
\[
\frac{\sqrt[3]{A}}{\sqrt[3]{B}}
=\sqrt[3]{\frac AB},\qquad B\ne0.
\]
Такой путь позволяет избежать путаницы при работе с несколькими
дробными показателями.
\subsection*{Исправление от последнего правильного шага}
Возьмём исходную дробь и сначала полностью преобразуем числитель:
\[
\bigl(\sqrt[3]{7x^6}\bigr)^2
=\sqrt[3]{(7x^6)^2}
=\sqrt[3]{49x^{12}}.
\]
Знаменатель пока оставим без изменений. Получается
\[
\frac{\sqrt[3]{49x^{12}}}{\sqrt[3]{49x^4}}.
\]
Теперь можно объединить оба кубических корня в один:
\[
\sqrt[3]{\frac{49x^{12}}{49x^4}}.
\]
Так как $x\ne0$, сокращение на $49x^4$ допустимо:
\[
\sqrt[3]{\frac{49x^{12}}{49x^4}}
=\sqrt[3]{x^{12-4}}=\sqrt[3]{x^8}.
\]
\subsection*{Полное правильное решение}
\textbf{Шаг 1.} Запишем выражение и отметим, что знаменатель
не равен нулю при $x\ne0$:
\[
\frac{\bigl(\sqrt[3]{7x^6}\bigr)^2}{\sqrt[3]{49x^4}}.
\]
\textbf{Шаг 2.} Используем правило степени корня в числителе:
\[
\bigl(\sqrt[3]{7x^6}\bigr)^2
=\sqrt[3]{(7x^6)^2}
=\sqrt[3]{49x^{12}}.
\]
Здесь $7^2=49$, а $(x^6)^2=x^{6\cdot2}=x^{12}$.

\textbf{Шаг 3.} Разделим кубические корни:
\[
\frac{\sqrt[3]{49x^{12}}}{\sqrt[3]{49x^4}}
=\sqrt[3]{\frac{49x^{12}}{49x^4}}.
\]
Это возможно, потому что корень нечётной степени определён
для всех действительных подкоренных выражений, а знаменатель
не равен нулю.

\textbf{Шаг 4.} Сократим общий числовой множитель и вычтем
показатели одинаковых степеней:
\[
\sqrt[3]{\frac{49x^{12}}{49x^4}}
=\sqrt[3]{\frac{x^{12}}{x^4}}
=\sqrt[3]{x^{12-4}}
=\sqrt[3]{x^8}.
\]
\textbf{Шаг 5.} Выделим под кубическим корнем полный куб:
\[
x^8=x^6x^2=(x^2)^3\,x^2.
\]
Следовательно,
\[
\sqrt[3]{x^8}
=\sqrt[3]{(x^2)^3x^2}
=x^2\sqrt[3]{x^2}.
\]
Равенство верно при любом действительном $x\ne0$:
кубический корень из $(x^2)^3$ равен именно $x^2$.
\subsection*{Проверка}
Проверим, что последний переход действительно сохраняет значение.
Возведём полученное выражение в куб:
\[
\bigl(x^2\sqrt[3]{x^2}\bigr)^3
=x^6\cdot x^2=x^8.
\]
Значит, оно совпадает с $\sqrt[3]{x^8}$.

Дополнительно проверим знак: исходный числитель --- квадрат,
поэтому неотрицателен; исходный знаменатель при $x\ne0$
положителен, так как $49x^4>0$. Полученное
$x^2\sqrt[3]{x^2}$ также положительно при $x\ne0$.
Противоречий нет.
\subsection*{Итог}
\textcolor{teal!60!black}{Итог:} преобразования нужно
выполнять для числителя и знаменателя по отдельности,
сохраняя их места в дроби. После объединения кубических
корней получается
\[
\boxed{x^2\sqrt[3]{x^2}},\qquad x\ne0.
\]
\subsection*{Задача на закрепление}
При $y\ne0$ упростите выражение:
\[
\frac{\bigl(\sqrt[3]{5y^6}\bigr)^2}{\sqrt[3]{25y^5}}.
\]

\vfill
\begin{center}
\small Лёвин Артём Александрович эксклюзивно для Марка Гукина
\end{center}
\end{document}